% concatenated from main__2_.tex
\documentclass[12pt]{article}

\usepackage[T1]{fontenc}
\usepackage[utf8]{inputenc}
\usepackage{lmodern}
\usepackage{microtype}
\usepackage[letterpaper,margin=1in]{geometry}
\usepackage{amsmath,amssymb,amsthm,mathtools}
\usepackage{hyperref}

\newtheorem{theorem}{Theorem}
\newtheorem{lemma}[theorem]{Lemma}
\newtheorem{proposition}[theorem]{Proposition}
\newtheorem{corollary}[theorem]{Corollary}
\theoremstyle{remark}
\newtheorem{remark}[theorem]{Remark}

\newcommand{\one}{\mathbf{1}}
\newcommand{\rev}{\mathrm{rev}}
\newcommand{\cL}{\mathcal{L}}
\newcommand{\cX}{\mathcal{X}}
\newcommand{\margin}{M}

\title{A Linear Lower Bound for Dominating Sets in $k$-Majority Tournaments}
\author{Jiangdong Ai\thanks{School of Mathematical Sciences and LPMC,
    Nankai University. {\tt jd@nankai.edu.cn}.
    Funded by the National Natural Science Foundation of China
    (No.\,12522117, No.\,12401456),
     the Fundamental and Interdisciplinary Disciplines Breakthrough Plan
    of the Ministry of Education of China (JYB2025XDXM207),
    and the Fundamental Research Funds for the Central Universities,
    Nankai University.}
     \hspace{2mm}
     Xiangjie Yi\thanks{School of Mathematical Sciences and LPMC,
    Nankai University. {\tt 2310184@mail.nankai.edu.cn}.}}
\date{}

\begin{document}
\maketitle

\begin{abstract}
A $k$-majority tournament on a finite vertex set is defined by $2k-1$ linear orders, with $u\to v$ when $u$ lies above $v$ in at least $k$ of
the orders.  Let $F(k)$ be the maximum, over all $k$-majority
tournaments, of the size of a minimum dominating set.  Alon,
Brightwell, Kierstead, Kostochka, and Winkler proved that $C_1k/\log k \leq F(k) \leq C_2k\log k$ for suitable positive constants $C_1$ and $C_2$.  In this paper, we prove the linear lower bound $F(k)\ge \left\lfloor\frac{k+1}{2}\right\rfloor $ for $k\ge 3$.

\end{abstract}


\section{Introduction}

Let $T=(V,E)$ be a tournament.  For $x,y\in V$, we write $x\to y$ when
$xy\in E$.  If $X,Y\subseteq V$, then $X\to Y$ means that for every
$y\in Y\setminus X$ there is an $x\in X$ such that $x\to y$; for a single vertex, we write $x\to Y$ in place of $\{x\}\to Y$.  A set
$D\subseteq V$ is a \emph{dominating set} if $D\to V$, and the minimum
size of such a set is denoted by $\gamma(T)$.  For a positive integer
$n$, let $[n]=\{1,\ldots,n\}$.

Let $P_1,\ldots,P_{2k-1}$ be linear orders on a finite set $V$.  They
\emph{realize} the tournament $T$ in which $u\to v$ if and only if $u$
lies above $v$ in at least $k$ of the orders.  A tournament realizable
in this way is called a \emph{$k$-majority tournament}.  Following
Alon, Brightwell, Kierstead, Kostochka, and Winkler~\cite{ABKKW}, define
\[
  F(k)=\max\{\gamma(T):T\text{ is a $k$-majority tournament}\}.
\]
The existence of this maximum is part of the result in~\cite{ABKKW}.

The realization of finite asymmetric relations---in particular, tournaments---as strict majority relations of profiles of linear orders goes back to McGarvey~\cite{McGarvey}.  Let $v(n)$ denote the least number of linear orders that suffices to realize every $n$-vertex
tournament.  Stearns~\cite{Stearns} proved that $v(n)=\Omega(n/\log n)$
and that $v(n)\le n+2$, and Erd\H{o}s and Moser~\cite{ErdosMoser}
improved the upper bound to $v(n)=O(n/\log n)$.  The older papers
themselves treat the more general realization problem for oriented preference patterns, in which ties are permitted.
Alon~\cite{AlonVoting} later studied realizations with a uniform quantitative lower bound on the majority supporting every prescribed arc.

Moon~\cite{Moon} gives standard graph-theoretic background on
tournaments, while Laslier~\cite{Laslier} surveys tournament solutions
and majority voting.  Large transitive subtournaments in majority
tournaments were studied by Milans, Schreiber, and
West~\cite{MilansSchreiberWest}; recent work of Shapira and
Yuster~\cite{ShapiraYuster} gives sharper Ramsey-type bounds for this
problem.  Coste, Flesch, Laison, McNicholas, and
Miyata~\cite{CosteEtAl} introduced weighted $k$-majority tournaments and
an approval-gap parameter related to dominating sets.

Alon et al.~\cite{ABKKW} proved that $F(2)=3$, that $F(3)\ge 4$, and
that there are positive absolute constants $c_1,c_2$ such that
\begin{equation}\label{eq:known-bounds}
  c_1\frac{k}{\log k}\le F(k)\le c_2 k\log k.
\end{equation}
Fidler~\cite{Fidler} later obtained substantially better upper bounds for small $k$, including $F(3)\le 12$. More recently, Bourneuf, Charbit, and Thomassé~\cite{BCT}, and
independently Charikar, Ramakrishnan, and Wang~\cite{CRW}, proved that every $(1/2-\varepsilon)$-majority digraph---in which $u\to v$ whenever
$u$ lies above $v$ in at least a $(1/2-\varepsilon)$-fraction of the orders---has a dominating set of size $O(\varepsilon^{-2})$.  Lowering the threshold in this way can only add arcs, but for a profile of
$2k-1$ orders the largest attainable minority fraction is
$(k-1)/(2k-1)=\frac12-\frac{1}{4k-2}$, so any
$0<\varepsilon<\frac{1}{4k-2}$ makes the corresponding digraph coincide with the strict-majority tournament.  Consequently these results yield
only $F(k)=O(k^2)$ for $k$-majority tournaments and do not improve the $O(k\log k)$ upper bound in~\eqref{eq:known-bounds}.

The lower-bound construction of Alon et al.~\cite{ABKKW} is
based on a closely related incidence-difference principle.
Its vertices are disjoint pairs $(A,B)$, and pairs of linear
orders indexed by the elements of the ground set realize the
margin $2\bigl(|A\cap B'|-|A'\cap B|\bigr)$.
Given $t$ prescribed vertices $(A_i,B_i)$, their construction
chooses $A$ to contain all the sets $B_i$ and then chooses $B$
so that $B\setminus A_i\neq\varnothing$ for every $i$.
Satisfying these simultaneous escape conditions requires a set
$B$ of logarithmic size and leads to a realization using
$\Theta(t\log t)$ orders.

Our construction retains the incidence-difference principle but
distributes the avoidance tasks among the coordinates of a
Cartesian product.  Coordinate $i$ is responsible only for the
$i$th prescribed vertex, while an unused row or column makes all
negative incidences vanish in that coordinate and preserves one
positive incidence.  Since the resulting local grid margin is
realized by four orders independently of the grid size, this
coordinatewise localization removes the logarithmic loss.

A tournament has \emph{property $S_t$} if, for every $t$-element set
$U\subseteq V$, there is a vertex $x\in V\setminus U$ such that
$x\to U$. Thus a tournament on more than $t$ vertices has property $S_t$ if and
only if it has no dominating set of size at most $t$, that is, if and only if $\gamma(T)\ge t+1$. 
Hence, to prove that $F(k) \ge t+1$,
it suffices to realize a tournament with property $S_t$ by $2k-1$ linear orders. Whether
such tournaments exist for every $t$ was asked by Sch\"utte and answered by
Erd\H{o}s~\cite{ErdosProblem}, who gave a probabilistic existence proof; see also Szekeres and
Szekeres~\cite{SzekeresSzekeres}. Our main result is the following explicit construction.

\begin{theorem}\label{thm:construction}
For every positive integer $t$, there is a tournament $T_t$ on
$\bigl(2(t+1)^2\bigr)^t
$ vertices that is realized by $4t+1$ linear orders and has property
$S_t$.  In particular, $T_t$ is a $(2t+1)$-majority tournament and
$\gamma(T_t)\ge t+1$.
\end{theorem}

Padding a realization by pairs of mutually reverse orders gives the
following consequence.

\begin{corollary}\label{cor:linear}
For every integer $k\ge 3$, $F(k)\ge \left\lfloor\frac{k+1}{2}\right\rfloor$.
Consequently,
$\liminf_{k\to\infty}\frac{F(k)}{k}\ge \frac12$.
\end{corollary}

For small $k$ the displayed estimate is not best possible.  Indeed, the
bound $F(3)\ge4$ from~\cite{ABKKW}, together with padding by reverse
pairs, gives $F(k)\ge4$ for every $k\ge3$, which is at least as strong as Corollary~\ref{cor:linear} for $k\le8$. The significance of
Corollary~\ref{cor:linear} is that it improves the general lower bound
in~\eqref{eq:known-bounds} from order $k/\log k$ to order $k$.

The construction is exact rather than probabilistic.  Each local state is a line of the $(t+1)\times(t+1)$ grid together with a marked point on it, and the difference of two point--line incidences gives the antisymmetric margin function that the four local orders realize.

\section{Margins and a lifting lemma}

For a linear order $L$ on a finite set $V$, define, for distinct
$x,y\in V$,
\[
  \sigma_L(x,y)=
  \begin{cases}
    1,&\text{if $x$ lies above $y$ in $L$},\\
   -1,&\text{if $y$ lies above $x$ in $L$}.
  \end{cases}
\]
If $\cL=(L_1,\ldots,L_m)$ is a profile of linear orders, its
\emph{margin} is
\[
  \margin_{\cL}(x,y)=\sum_{j=1}^m \sigma_{L_j}(x,y).
\]
An \emph{antisymmetric margin function} on a finite set $V$ is an
integer-valued function $K:V\times V\to\mathbb Z$ satisfying
$K(x,x)=0$ and $K(x,y)=-K(y,x)$ for all $x,y\in V$.  We say that
$\cL$ \emph{realizes $K$ exactly} when $\margin_{\cL}(x,y)=K(x,y)$
for every pair of distinct vertices.

The following observation permits a local profile to be lifted through
a map without creating margins inside its fibers.

\begin{lemma}\label{lem:lifting}
Let $Z$ and $V$ be finite sets, let $\pi:V\to Z$, and suppose that
$L_1,\ldots,L_{2m}$ realize an antisymmetric margin function $K$ on
$Z$ exactly.  There are $2m$ linear orders on $V$ whose margin on
distinct $x,y\in V$ is
\[
  \widetilde K(x,y)=
  \begin{cases}
    K(\pi(x),\pi(y)),&\text{if $\pi(x)\ne\pi(y)$},\\
    0,&\text{if $\pi(x)=\pi(y)$}.
  \end{cases}
\]
\end{lemma}

\begin{proof}
Fix a linear order $R$ on $V$.  For each $j\in[2m]$, order the fibers
$\pi^{-1}(z)$ according to the order $L_j$ of their labels.  Inside
each fiber use the restriction of $R$ in the first $m$ orders and the
restriction of $R^{\rev}$ in the last $m$ orders.

If $\pi(x)\ne\pi(y)$, the relative position of the two fibers is
prescribed by $L_j$, independently of the order inside either fiber.
The resulting margin is therefore $K(\pi(x),\pi(y))$.  If
$\pi(x)=\pi(y)$, the first $m$ contributions are cancelled by the last
$m$ contributions.
\end{proof}

Margins add when profiles are concatenated.  We shall therefore apply
Lemma~\ref{lem:lifting} independently in several coordinates and sum
the resulting margin functions.

\section{A four-order grid construction}

Fix an integer $q\ge2$, and let
$\Gamma_q=[q]\times[q]$
be the $q\times q$ grid.  For $r,a\in[q]$, let $H_{r,a}$ be the
horizontal state with supporting line and marked point
\[
  \ell(H_{r,a})=\{(u,r):u\in[q]\},
  \qquad
  p(H_{r,a})=(a,r).
\]
For $c,d\in[q]$, let $V_{c,d}$ be the vertical state with
\[
  \ell(V_{c,d})=\{(c,v):v\in[q]\},
  \qquad
  p(V_{c,d})=(c,d).
\]
Let
\[
  Z_q=\{H_{r,a}:r,a\in[q]\}\mathbin{\dot\cup}
      \{V_{c,d}:c,d\in[q]\}.
\]
Thus $|Z_q|=2q^2$, and the marked point of every state lies on its
supporting line.

Define the antisymmetric margin function $\psi$ on $Z_q$ by
\begin{equation}\label{eq:psi}
  \psi(z,w)=
  \one_{\{p(z)\in\ell(w)\}}-
  \one_{\{p(w)\in\ell(z)\}}.
\end{equation}
The same formula gives $\psi(z,z)=0$.

\begin{lemma}\label{lem:psi-values}
For all admissible indices,
\[
  \psi(H_{r,a},H_{s,b})=0,
  \qquad
  \psi(V_{c,d},V_{e,f})=0,
\]
and
\begin{equation}\label{eq:cross-value}
  \psi(H_{r,a},V_{c,d})=
  \one_{\{a=c\}}-\one_{\{r=d\}}.
\end{equation}
\end{lemma}

\begin{proof}
For two horizontal states, each incidence in~\eqref{eq:psi} is
equivalent to equality of the two row indices, so the two terms cancel.
The vertical--vertical case is the same with rows replaced by columns.
Finally, $(a,r)$ lies on the supporting column of $V_{c,d}$ exactly
when $a=c$, while $(c,d)$ lies on the supporting row of $H_{r,a}$
exactly when $r=d$.  This gives~\eqref{eq:cross-value}.
\end{proof}

\begin{proposition}\label{prop:four-orders}
The margin function $2\psi$ on $Z_q$ is realized exactly by four linear
orders.
\end{proposition}

\begin{proof}
For each $a\in[q]$, form the ordered blocks
\[
  \mathbf H^a=(H_{1,a},\ldots,H_{q,a}),
  \qquad
  \mathbf V^a=(V_{a,1},\ldots,V_{a,q}).
\]
We write blocks in succession to mean that every element of an earlier block lies above every element of a later one. Define the two column orders
\begin{align*}
  C_1&=\mathbf H^1\mathbf V^1\,
       \mathbf H^2\mathbf V^2\cdots
       \mathbf H^q\mathbf V^q,\\
  C_2&=(\mathbf H^q)^{\rev}(\mathbf V^q)^{\rev}\cdots
       (\mathbf H^1)^{\rev}(\mathbf V^1)^{\rev}.
\end{align*}
Note that $C_2$ is not the reverse of $C_1$: the pairs $\mathbf H^a\mathbf V^a$ occur in the reverse order and each block is
reversed internally, but $\mathbf H^a$ still precedes $\mathbf V^a$ for every $a$.

Consider first two states of the same type, say $H_{r,a}$ and
$H_{s,b}$.  If $a=b$, both lie in the block $\mathbf H^a$, whose
internal order is reversed in $C_2$; if $a\ne b$, the blocks
$\mathbf H^a$ and $\mathbf H^b$ occur in opposite orders in $C_1$ and $C_2$.  Either way the two contributions cancel, and the same argument
applies to vertical--vertical pairs. Hence the column pair has zero
margin on pairs of the same type.

Now consider $H_{r,a}$ and $V_{c,d}$.  If $a=c$, then $H_{r,a}$ lies above $V_{c,d}$ in both column orders, by the remark above.  If
$a\ne c$, the blocks $\mathbf H^a$ and $\mathbf V^c$ occur in opposite orders in $C_1$ and $C_2$.  Therefore
\begin{equation}\label{eq:column-pair}
  \sigma_{C_1}(H_{r,a},V_{c,d})+
  \sigma_{C_2}(H_{r,a},V_{c,d})
  =2\one_{\{a=c\}}.
\end{equation}

For each $r\in[q]$, form the blocks
\[
  \widehat{\mathbf V}^{\,r}=(V_{1,r},\ldots,V_{q,r}),
  \qquad
  \widehat{\mathbf H}^{\,r}=(H_{r,1},\ldots,H_{r,q}),
\]
and define the two row orders
\begin{align*}
  R_1&=\widehat{\mathbf V}^{\,1}\widehat{\mathbf H}^{\,1}\,
       \widehat{\mathbf V}^{\,2}\widehat{\mathbf H}^{\,2}\cdots
       \widehat{\mathbf V}^{\,q}\widehat{\mathbf H}^{\,q},\\
  R_2&=(\widehat{\mathbf V}^{\,q})^{\rev}
       (\widehat{\mathbf H}^{\,q})^{\rev}\cdots
       (\widehat{\mathbf V}^{\,1})^{\rev}
       (\widehat{\mathbf H}^{\,1})^{\rev}.
\end{align*}
Again $R_2$ is not the reverse of $R_1$: the pairs
$\widehat{\mathbf V}^{\,r}\widehat{\mathbf H}^{\,r}$ occur in the
reverse order and each block is reversed internally, but
$\widehat{\mathbf V}^{\,r}$ still precedes $\widehat{\mathbf H}^{\,r}$
for every $r$.  Exactly as above, the row pair has zero margin on pairs
of the same type.  For a cross pair, if $r=d$ then $V_{c,d}$ lies above
$H_{r,a}$ in both row orders, while if $r\ne d$ the blocks
$\widehat{\mathbf H}^{\,r}$ and $\widehat{\mathbf V}^{\,d}$ occur in
opposite orders.  Therefore
\begin{equation}\label{eq:row-pair}
  \sigma_{R_1}(H_{r,a},V_{c,d})+
  \sigma_{R_2}(H_{r,a},V_{c,d})
  =-2\one_{\{r=d\}}.
\end{equation}
Equations~\eqref{eq:column-pair} and~\eqref{eq:row-pair}, together with
Lemma~\ref{lem:psi-values}, show that $C_1,C_2,R_1,R_2$ have total
margin $2\psi$ on every pair of distinct states.
\end{proof}

\section{The product construction}

Fix a positive integer $t$ and let $q=t+1$.  The vertex set of the
global construction is
$\cX_t=Z_q^t$.
Write $x=(x_1,\ldots,x_t)$ for a vertex of $\cX_t$, and define
\begin{equation}\label{eq:Phi}
  \Phi(x,y)=\sum_{g=1}^t \psi(x_g,y_g).
\end{equation}
The function $\Phi$ is antisymmetric.

\begin{proposition}\label{prop:global-orders}
There are $4t$ linear orders on $\cX_t$ whose margin is $2\Phi$.
\end{proposition}

\begin{proof}
For $g\in[t]$, let $\pi_g:\cX_t\to Z_q$ be the projection
$\pi_g(x)=x_g$.  Apply Lemma~\ref{lem:lifting} to the four orders in
Proposition~\ref{prop:four-orders} and the map $\pi_g$.  The resulting
four orders on $\cX_t$ have margin $2\psi(x_g,y_g)$ on $x,y$.
Concatenating these profiles for all $g\in[t]$ gives $4t$ orders with
margin
\[
  2\sum_{g=1}^t\psi(x_g,y_g)=2\Phi(x,y).
\]
\end{proof}

The next lemma is the point at which the choice $q=t+1$ is used.

\begin{lemma}\label{lem:common-vertex}
For any distinct $y^1,\ldots,y^t\in\cX_t$, there is a vertex
$x\in\cX_t\setminus\{y^1,\ldots,y^t\}$ such that
$\Phi(x,y^i)\ge1$ for every $i\in[t]$.
\end{lemma}

\begin{proof}
We choose $x=(x_1,\ldots,x_t)$ one coordinate at a time.  Coordinate
$g$ will give a positive contribution against $y^g$ and a nonnegative
contribution against every selected vertex.

Fix $g\in[t]$ and consider the set of marked points
\[
  P_g=\{p(y^1_g),\ldots,p(y^t_g)\}\subseteq\Gamma_q.
\]
Suppose first that $y^g_g=H_{r,a}$.  The points in $P_g$ occupy at most
$t$ columns, whereas $\Gamma_q$ has $q=t+1$ columns.  Choose a column
$c$ containing no point of $P_g$, and let
$x_g=V_{c,r}$.
The supporting line of $x_g$ contains none of the marked points
$p(y^i_g)$.  Hence~\eqref{eq:psi} gives
\[
  \psi(x_g,y^i_g)=
  \one_{\{p(x_g)\in\ell(y^i_g)\}}\ge0
  \qquad (i\in[t]).
\]
Since $p(x_g)=(c,r)$ lies on the supporting row of $y^g_g$, we also
have $\psi(x_g,y^g_g)=1$.

Suppose instead that $y^g_g=V_{c,d}$.  The points in $P_g$ occupy at
most $t$ rows.  Choose a row $r$ containing no point of $P_g$, and let $x_g=H_{r,c}$.
The same argument, with rows and columns interchanged, gives
\[
  \psi(x_g,y^i_g)\ge0\quad(i\in[t]),
  \qquad
  \psi(x_g,y^g_g)=1.
\]

Now fix $i\in[t]$.  Every term in~\eqref{eq:Phi} is nonnegative, and
the term with $g=i$ is one.  Therefore $\Phi(x,y^i)\ge1$.  Moreover,
$x_i$ and $y^i_i$ have opposite types, so $x_i\ne y^i_i$.  Hence
$x\ne y^i$ for every $i$, and $x$ lies outside the selected set.
\end{proof}

\section{Proof of the main results}

\begin{proof}[Proof of Theorem~\ref{thm:construction}]
By Proposition~\ref{prop:global-orders}, there are $4t$ linear orders
on $\cX_t$ with margin $2\Phi$.  Add an arbitrary linear order $L_*$ on
$\cX_t$.  The margin of the resulting odd profile is
\begin{equation}\label{eq:final-margin}
  \margin(x,y)=2\Phi(x,y)+\sigma_{L_*}(x,y).
\end{equation}
If $\Phi(x,y)\ge1$, then $\margin(x,y)\ge1$; if $\Phi(x,y)\le-1$, then
$\margin(x,y)\le-1$.  When $\Phi(x,y)=0$, the last order breaks the
tie.  Thus the profile realizes a tournament, and every pair having
nonzero $\Phi$-score is oriented according to the sign of $\Phi$.

The profile has $4t+1=2(2t+1)-1$ orders, so the tournament is a $(2t+1)$-majority tournament.  Given any
$t$ vertices $y^1,\ldots,y^t$, Lemma~\ref{lem:common-vertex} supplies
an outside vertex $x$ with $\Phi(x,y^i)\ge1$ for every $i$.  By
\eqref{eq:final-margin}, $x\to y^i$ for every $i$, and hence the
tournament has property $S_t$.

Finally, suppose that $D$ is a dominating set with $|D|\le t$.  Extend
$D$, if necessary, to a $t$-element set $U$.  Property $S_t$ gives a
vertex $x\notin U$ with $x\to U$, so no vertex of $D$ dominates $x$.
This contradiction shows that $\gamma(T_t)>t$, and therefore
$\gamma(T_t)\ge t+1$.
\end{proof}

\begin{proof}[Proof of Corollary~\ref{cor:linear}]
Let $k\ge3$, and let
$t=\left\lfloor\frac{k-1}{2}\right\rfloor$ and 
 $K=2t+1$.
Then $K\le k$.  The tournament $T_t$ in Theorem~\ref{thm:construction}
is realized by $2K-1$ orders.  If $K<k$, append $k-K$ pairs
$(L,L^{\rev})$ of mutually reverse linear orders.  Every such pair has
zero margin on every pair of vertices, so the tournament is unchanged,
while the number of orders becomes
\[
  (2K-1)+2(k-K)=2k-1.
\]
Thus $T_t$ is also a $k$-majority tournament, and
\[
  F(k)\ge\gamma(T_t)\ge t+1
  =\left\lfloor\frac{k+1}{2}\right\rfloor.
\]
Dividing by $k$ and taking the lower limit proves the final assertion.
\end{proof}

\section{Concluding remarks}

The construction uses four local orders for each vertex protected by
property $S_t$.  Its two essential features are the exact realization
of the marked row--column margin function and the availability of one
more row and column than the number of prescribed vertices.  No probabilistic
selection or asymptotic compression is used.

The constant $1/2$ has not been optimized.  One possible direction is
to replace the row--column gadget by a local incidence system that
protects several prescribed vertices per coordinate without a
proportional increase in the number of orders.  In view of the upper
bound in~\eqref{eq:known-bounds}, the principal remaining asymptotic
question is whether $F(k)=O(k)$.

\section*{Declaration on the use of generative AI}

During the preparation of this manuscript, the authors used OpenAI's
ChatGPT (5.6 Sol) in a limited supporting role for language polishing, presentation, bibliographic checks, and exploratory mathematical discussion.  Some mathematical ideas and elements of the construction
arose through interaction between the authors and these tools.  The authors directed the mathematical investigation, selected and developed the final construction and arguments, verified all proofs and claims, and take full responsibility for the content of the manuscript.
\begin{thebibliography}{99}

\bibitem{ABKKW}
N.~Alon, G.~Brightwell, H.~A.~Kierstead, A.~V.~Kostochka and P.~Winkler,
Dominating sets in $k$-majority tournaments,
\emph{J. Combin. Theory Ser. B} 96 (2006), 374--387.

\bibitem{AlonVoting}
N.~Alon,
Voting paradoxes and digraphs realizations,
\emph{Adv. in Appl. Math.} 29 (2002), 126--135.
\bibitem{BCT}
R.~Bourneuf, P.~Charbit, and S.~Thomassé,
A dense neighborhood lemma: Applications of partial concept
classes to domination and chromatic number,
in: \emph{2025 IEEE 66th Annual Symposium on Foundations of
Computer Science (FOCS)}, 2025.

\bibitem{CRW}
M.~Charikar, P.~Ramakrishnan, and K.~Wang,
Approximately dominating sets in elections,
in: \emph{Proceedings of the 2026 Annual ACM--SIAM Symposium
on Discrete Algorithms (SODA)}, SIAM, 2026, 1747--1760.

\bibitem{CosteEtAl}
J.~Coste, B.~Flesch, J.~D.~Laison, E.~McNicholas and D.~Miyata,
Approval gap of weighted $k$-majority tournaments,
\emph{Theory Appl. Graphs} 11 (2024), no.~1, Article~3, 1--14.

\bibitem{ErdosProblem}
P.~Erd\H{o}s,
On a problem in graph theory,
\emph{Math. Gaz.} 47 (1963), 220--223.

\bibitem{ErdosMoser}
P.~Erd\H{o}s and L.~Moser,
On the representation of directed graphs as unions of orderings,
\emph{Magyar Tud. Akad. Mat. Kutat\'o Int. K\"ozl.} 9 (1964), 125--132.

\bibitem{Fidler}
D.~Fidler,
A recurrence for bounds on dominating sets in $k$-majority tournaments,
\emph{Electron. J. Combin.} 18 (2011), Paper~P166.

\bibitem{Laslier}
J.-F.~Laslier,
\emph{Tournament Solutions and Majority Voting},
Springer, Berlin, 1997.

\bibitem{McGarvey}
D.~C.~McGarvey,
A theorem on the construction of voting paradoxes,
\emph{Econometrica} 21 (1953), 608--610.

\bibitem{MilansSchreiberWest}
K.~G.~Milans, D.~H.~Schreiber and D.~B.~West,
Acyclic sets in $k$-majority tournaments,
\emph{Electron. J. Combin.} 18 (2011), Paper~P122.

\bibitem{Moon}
J.~W.~Moon,
\emph{Topics on Tournaments},
Holt, Rinehart and Winston, New York, 1968.
\bibitem{ShapiraYuster}
A.~Shapira and R.~Yuster,
On Ramsey properties of $k$-majority tournaments,
\emph{arXiv:2603.04174} (2026).
\bibitem{Stearns}
R.~Stearns,
The voting problem,
\emph{Amer. Math. Monthly} 66 (1959), 761--763.
\bibitem{SzekeresSzekeres}
E.~Szekeres and G.~Szekeres,
On a problem of Sch\"utte and Erd\H{o}s,
\emph{Math. Gaz.} 49 (1965), no.~369, 290--293.
\end{thebibliography}

\end{document}
